\input{../KU-Beamer-Template/style/template.tex}
\usepackage{../KU-Beamer-Template/style/koc} 
\usepackage{minted}
\usepackage{upquote}
 
\title{KOLT Python} 
\subtitle{Basic Operators \& Branching} 
\newdate{date}{25}{02}{2019}
\date{\displaydate{date}}
\author{Ahmet Uysal}

\titlegraphic{\includegraphics[scale=0.2]{../KU-Beamer-Template/style/images/logo_kolt.eps}}

\setbeamercovered{invisible} % transparent

\begin{document}
    \maketitle
    \frame{\frametitle{Agenda}\tableofcontents}

    \section{Recap}

        \begin{frame}{Comments}
            \LARGE
            \inputminted[frame=single,framesep=2pt]{python3}{../Lecture1/code-examples/comments.py}
            Python will basically ignore comments, they are purely written \textbf{for humans}!
        \end{frame}

        \begin{frame}{Variables}
            \LARGE
            \begin{table}[]
                \resizebox{0.8\textwidth}{!}{
                \begin{tabular}{|l|l|l|}
                \hline
                \textcolor{koc}{\textbf{Type}} &\textcolor{koc}{\textbf{Explanation}} & \textcolor{koc}{\textbf{Examples}} \\ \hline
                \textbf{\texttt{int}}  & represent \textbf{integers} & 3, 4, 17, -10 \\ \hline 
                \textbf{\texttt{float}} & represent \textbf{real numbers} & 3.0, 1.11, -109.123123 \\ \hline
                \textbf{\texttt{bool}} & represent \textbf{boolean} truth values & \texttt{True}, \texttt{False} \\ \hline
                \textbf{\texttt{str}} & A sequence of characters. & \textquotesingle Hello\textquotesingle, \textquotesingle \textquotesingle, \textquotesingle 3\textquotesingle \\ \hline
                \textbf{\texttt{NoneType}}& special and has one value, None & \texttt{None} \\ \hline
                \end{tabular}}
            \end{table}

            \begin{itemize}
                \item How to create a variable?        
                \texttt{variable\_name = value}
                \item How about types?
                use \texttt{type()}
                \item Can a variable change type?
                \textbf{Yes!} Just assing a new value with any type.
                \item What if we if want to convert a value between types, i.e, \textquotesingle 2\textquotesingle $\to$ 2
            \end{itemize}
        \end{frame}

        \begin{frame}{Casting}
            \LARGE
            \begin{itemize}
                \item \texttt{int(\textquotesingle 2\textquotesingle)} $\to$ 2
                \item Any possible reasons for casting?
                -taking user input
                -reading numbers from a file?
                \item Can we cast every value to every type?
                \textbf{NO!} try \texttt{int(\textquotesingle hello\textquotesingle)}
            \end{itemize}
        \end{frame}

        \begin{frame}{Console I/O(Input/Output)}
            \huge
            \textbf{\texttt{print(*args, sep=\textquotesingle \ \textquotesingle, end=\textquotesingle \textbackslash n\textquotesingle )}}
            \begin{itemize}
                \LARGE
                \item Can take arbitrary number of arguments
                \item Separates elements with space by default
                \item Adds newline character \texttt{\textquotesingle \textbackslash n\textquotesingle} to end by default
            \end{itemize}
            
            \textbf{\texttt{input([prompt])}}
            \begin{itemize}
                \LARGE
                \item Prints the prompt to Console
                \item Program is paused until user enters something
                \item \textbf{returns an \texttt{str} object!} 
            \end{itemize}
        \end{frame}

\end{document}